\chapter{Fourier-Stieltjes transforms (one variable)}
\label{ch:III}

\section{Basic Properties}
\label{sec:III-1}

We assume as known the fundamentals of the theory of Riemann--Stieltjes integrals.

Let $F(y)$ be a function of bounded variation for $-\infty<y<\infty$. For $x$ real, let
\begin{equation}\tag{1.1}\label{eq:III-1.1}
\varphi(x)=\int_{-\infty}^{\infty}e^{ixy}\,dF(y)\equiv\lim_{R\to\infty}\int_{-R}^{R}e^{ixy}\,dF(y),\qquad(R>0).
\end{equation}
We call $\varphi$ the Fourier--Stieltjes transform of $F$, or the Fourier transform of $dF$, and denote it sometimes by the symbol $\hat dF$.

If, in particular,
\begin{equation}\tag{1.2}\label{eq:III-1.2}
F(y)=\int_{-\infty}^{y}f(t)\,dt,\qquad f\in L_{1}(-\infty,\infty),
\end{equation}
then \eqref{eq:III-1.1} reduces to the Fourier transform on $L_{1}(-\infty,\infty)$ studied in Ch.~I.

The integral in \eqref{eq:III-1.1} converges absolutely and uniformly and $\varphi(x)$ is a bounded, continuous function of $x$ defined for every $x$ in $(-\infty,\infty)$. We have only to note that
\[
\bigl|\varphi(x)\bigr|\le\int_{-\infty}^{\infty}\bigl|dF(y)\bigr|<\infty,
\]
and for any real $h\neq0$,
\begin{align*}
&\bigl|\varphi(x+h)-\varphi(x)\bigr|
\le\int_{-\infty}^{\infty}\bigl|e^{ihy}-1\bigr|\,\bigl|dF(y)\bigr|\\
&\le|h|\int_{|y|<R}|y|\,\bigl|dF(y)\bigr|+2\int_{|y|\ge R}\bigl|dF(y)\bigr|=I_{1}+I_{2},
\end{align*}
say, where $R>0$. Given any $\varepsilon>0$, one can choose $R$ so large that $|I_{2}|<\varepsilon/2$, and $h$ so small that $|I_{1}|<\varepsilon/2$.

But unlike the Fourier transform on $L_{1}(-\infty,\infty)$, $\varphi(x)$ does not necessarily tend to zero as $|x|\to\infty$. For example, if $F(x)=1$, for $x>0$; $F(x)=-1$, for $x<0$; and $F(0)=0$, then $\varphi(x)=2$.

\begin{theorem}\label{thm:III-1}
Let $F(x)$ be of bounded variation in $(-\infty,\infty)$, with
\begin{equation}\tag{1.3}\label{eq:III-1.3}
F(x)=\tfrac12\bigl\{F(x+0)+F(x-0)\bigr\},\qquad\text{for all }x,
\end{equation}
and
\begin{equation}\tag{1.4}\label{eq:III-1.4}
\varphi(x)=\int_{-\infty}^{\infty}e^{ixy}\,dF(y).
\end{equation}
Then we have
\begin{equation}\tag{1.5}\label{eq:III-1.5}
F(x)-F(0)=\frac{1}{2\pi}\lim_{R\to\infty}\int_{-R}^{R}\varphi(t)\,\frac{e^{-itx}}{-it}\,dt
\equiv\frac{1}{2\pi}\int_{-\infty}^{\infty}\varphi(t)\,\frac{e^{-itx}}{-it}\,dt,
\end{equation}
so that $\varphi$ determines $F$ up to an arbitrary, additive constant.
\end{theorem}

\begin{proof}
If $h$ is real, and fixed, and $h\neq0$, then \eqref{eq:III-1.4} gives
\begin{equation}\tag{1.6}\label{eq:III-1.6}
\varphi(x)\,e^{-ihx}=\int_{-\infty}^{\infty}e^{ix(y-h)}\,dF(y)=\int_{-\infty}^{\infty}e^{ixy}\,dF(y+h).
\end{equation}
From this and \eqref{eq:III-1.4} we get
\begin{equation}\tag{1.7}\label{eq:III-1.7}
\varphi(x)\bigl[e^{-ihx}-1\bigr]=\int_{-\infty}^{\infty}e^{ixy}\,dG(y),\qquad G(y)=F(y+h)-F(y).
\end{equation}

Here $G$ is of bounded variation in $(-\infty,\infty)$, and $G\in L_{1}(-\infty,\infty)$. For if we suppose $h\ge0$, $h$ fixed; and suppose that $F$ is non-decreasing (since $F$ is expressible as the difference of two non-decreasing, bounded functions, if it is real-valued; and if it is complex-valued, one can consider the real and imaginary parts separately), then $G(y)\ge0$, and we have for $R>0$,
\[
\begin{aligned}
\int_{-R}^{R}G(x)\,dx &= \int_{-R}^{R}\bigl\{F(x+h)-F(x)\bigr\}\,dx\\
&= \int_{R}^{R+h}F(x)\,dx-\int_{-R}^{-R+h}F(x)\,dx\\
&\to h\bigl\{F(+\infty)-F(-\infty)\bigr\},
\end{aligned}
\]
as $R\to+\infty$, and the limit is finite by hypothesis.

We note that $G(x)\to0$, as $|x|\to\infty$ (by definition and by the hypothesis on $F$), and by partial integration,
\[
\begin{aligned}
&\lim_{R\to\infty}  \Bigl[\int_{-R}^{R}e^{ixy}\,dG(y)\Bigr]\\
&= \lim_{R\to\infty}  \Bigl[e^{ixy}G(y)\Bigr]_{y=-R}^{y=R} - ix\lim_{R\to\infty}  \int_{-R}^{R}e^{ixy}G(y)\,dy,
\end{aligned}
\]
which gives
\[
\int_{-\infty}^{\infty}e^{ixy}\,dG(y)=-ix\int_{-\infty}^{\infty}e^{ixy}G(y)\,dy,
\]
so that, by \eqref{eq:III-1.7}, we have
\[
\varphi(x)\begin{bmatrix}\dfrac{e^{-ihx}-1}{-ix}\end{bmatrix}
=\int_{-\infty}^{\infty}e^{ixy}G(y)\,dy.
\]
Since $F(x)=\frac12[F(x+0)+F(x-0)]$, this implies (by Theorem~5, Ch.~I) that
\[
G(x)=\frac{1}{2\pi}\lim_{R\to\infty}\int_{-R}^{R}\varphi(y)\,\frac{e^{-ihy}-1}{-iy}\,e^{-ixy}\,dy,
\]
for each $x$. For $x=0$, this gives \eqref{eq:III-1.5}.
\end{proof}

\begin{theorem}\label{thm:III-2}
Let $F$ be of bounded variation in $(-\infty,+\infty)$, and let $F(x)=\frac12[F(x+0)+F(x-0)]$, for all $x$, and let $\varphi(x)=\int_{-\infty}^{\infty}e^{ixy}\,dF(y)$. Then we have
\begin{equation}\tag{1.8}\label{eq:III-1.8}
\int_{0}^{x}\bigl\{F(y)-F(-y)\bigr\}\,dy=\frac{1}{\pi}\int_{-\infty}^{\infty}\varphi(y)\,\frac{1-\cos xy}{y^{2}}\,dy,
\end{equation}
the integral on the right-hand side converging absolutely (since $\varphi$ is bounded).
\end{theorem}

\begin{proof}
We have, for $R>0$,
\[
\frac{1}{R\pi}\int_{-\infty}^{\infty}\varphi(y)\,\frac{1-\cos Ry}{2}\,dy
=\frac{1}{\pi}\int_{-\infty}^{\infty}dF(y)\int_{-\infty}^{\infty}e^{ixy}\,\frac{1-\cos Ry}{2}\,dy
\]
\[
=\int_{-R}^{R}\Bigl(1-\frac{|x|}{R}\Bigr)\,dF(x)
=\frac{1}{R}\int_{0}^{R}\bigl[F(y)-F(-y)\bigr]\,dy.
\]
(see Ex.2, \S\ref{sec:I-1}, Ch.~I; also Ex.1, \S\ref{sec:I-8}, Ch.~I).
\end{proof}

\subsection*{Remarks}

Let $F_{1}(x)$, $F_{2}(x)$, \dots, be non-decreasing, and bounded functions in $(-\infty<x<\infty)$, and let
\begin{equation}\tag{1.9}\label{eq:III-1.9}
\varphi_{n}(x)=\int_{-\infty}^{\infty}e^{ixy}\,dF_{n}(y),\qquad n=1,2,\dots.
\end{equation}
The following examples show the difficulty in preserving the equality sign in \eqref{eq:III-1.9} after letting $n\to+\infty$.

\noindent\textbf{(1.10)} Let $F_{n}(x)=0$, for $x<n$; $F_{n}(x)=1$, for $x\ge n$. Then $F_{n}(x)\to0$, for each $x$, as $n\to+\infty$, but $\varphi_{n}(x)=e^{inx}$ does not tend to a limit as $n\to+\infty$, except for $x=2\pi k$, where $k$ is an integer.

\noindent\textbf{(1.11)} Let $F_{n}(x)$ be continuous, $F_{n}(x)=0$ for $x\le-n$, $F_{n}(x)=1$ for $x\ge n$, $F_{n}(x)=\frac{x}{2n}+\frac12$ for $-n<x<n$ (linear). Then $F_{n}(x)\to\frac12$, for each $x$, as $n\to\infty$, and
\[
\varphi_{n}(x)=\int_{-n}^{n}\frac{e^{ixy}}{2n}\,dy=\frac{\sin nx}{nx},
\]
so that $\varphi_{n}(x)\to0$, as $n\to\infty$, for each $x\neq0$, and $\varphi_{n}(x)\to1$, at $x=0$, as $n\to\infty$, so that
\[
\lim_{n\to\infty}\varphi_{n}(x)\neq\int_{-\infty}^{\infty}e^{ixy}\,d\Bigl(\lim_{n\to\infty}F_{n}(y)\Bigr).
\]

\section{Distribution functions, and characteristic functions}
\label{sec:III-2}

In order to avoid some of the difficulties pointed up by the above examples, we shall now consider non-decreasing functions $F(x)$, which have finite limits at $x=\pm\infty$. We suppose further that
\begin{equation}\tag{2.1}\label{eq:III-2.1}
F(-\infty)=0,\qquad F(+\infty)=1.
\end{equation}
We call such a function $F$ a distribution function. We call the Fourier--Stieltjes transform of such an $F$, namely
\[
\varphi(x)=\int_{-\infty}^{\infty}e^{ixy}\,dF(y),
\]
a characteristic function, corresponding to the given distribution function $F$.

It is clear that $\varphi$ is continuous and bounded. (Note that $|\varphi(x)|\le F(+\infty)-F(-\infty)=1$, and $\varphi(0)=1$). It is also hermitian, in the sense that $\overline{\varphi(-x)}=\varphi(x)$.

Theorem~\ref{thm:III-1} implies that given $\varphi$, the distribution function $F$ is defined uniquely at the points of continuity of $F$, because $F(-\infty)=0$, $F(+\infty)=1$. The next theorem shows how the convergence of a sequence of distribution functions implies the convergence of the corresponding sequence of characteristic functions.

\begin{theorem}\label{thm:III-3}
Let $(F_{n})$ be a sequence of distribution functions, and $(\varphi_{n})$ the corresponding sequence of characteristic functions, so that
\begin{equation}\tag{2.2}\label{eq:III-2.2}
\varphi_{n}(x)=\int_{-\infty}^{\infty}e^{ixy}\,dF_{n}(y).
\end{equation}
If $F_{n}$ converges to a distribution function $F$ as $n\to\infty$, at the points of continuity of $F$, and if $\varphi$ is the characteristic function of $F$, then
\begin{equation}\tag{2.3}\label{eq:III-2.3}
\varphi_{n}(x)\to\varphi(x)\qquad\text{as }n\to\infty,
\end{equation}
and the convergence is uniform in every finite interval.
\end{theorem}

\begin{proof}
We note that the set of discontinuities of $F$ is countable. Given $\varepsilon>0$, we choose $R>0$, such that $F(x)$ is continuous both at $x=R$ and $x=-R$, with $F(-R)<\varepsilon$, $F(R)>1-\varepsilon$, and such that $F_{n}(\pm R)\to F(\pm R)$, as $n\to\infty$. Then we have, for $n\ge n_{0}$, also
\[
F_{n}(-R)<\varepsilon,\qquad F_{n}(R)>1-\varepsilon.
\]
Now
\begin{align*}
&\varphi(x)-\varphi_{n}(x)
=\int_{|y|\le R}e^{ixy}\,d\bigl\{F(y)-F_{n}(y)\bigr\}+\int_{|y|>R}e^{ixy}\,dF(y)\\
&\qquad-\int_{|y|>R}e^{ixy}\,dF_{n}(y)=I_{1}+I_{2}+I_{3},\quad\text{say.}
\end{align*}
Then we have
\[
|I_{2}|\le\bigl(1-F(R)+F(-R)\bigr)<2\varepsilon,\qquad\text{for all }x,
\]
while
\[
|I_{3}|\le\bigl(1-F_{n}(R)+F_{n}(-R)\bigr)<2\varepsilon,\qquad\text{for }n\ge n_{0},\text{ and all }x.
\]
On the other hand,
\[
I_{1}=\Bigl[\bigl(F(y)-F_{n}(y)\bigr)e^{ixy}\Bigr]_{y=-R}^{y=+R}-ix\int_{|y|\le R}\bigl\{F(y)-F_{n}(y)\bigr\}e^{ixy}\,dy.
\]
Since $F_{n}(\pm R)\to F(\pm R)$, as $n\to\infty$, the first term on the right-hand side tends to zero, uniformly in $x$. And
\[
\Bigl|-ix\int_{|y|\le R}\bigl\{F(y)-F_{n}(y)\bigr\}e^{ixy}\,dy\Bigr|
\le|x|\int_{|y|\le R}\bigl|F(y)-F_{n}(y)\bigr|\,dy\to0,
\]
uniformly in each finite $x$-interval since $|F|\le1$, $|F_{n}|\le1$; hence so does $I_{1}$. Thus we have altogether
\[
\bigl|\varphi(x)-\varphi_{n}(x)\bigr|<\varepsilon+2\varepsilon+2\varepsilon,\qquad\text{for }n\ge n_{0},
\]
and for $x$ in any finite interval.
\end{proof}

The next result is a kind of converse to Theorem~\ref{thm:III-3}.

\begin{theorem}\label{thm:III-4}
If $\varphi_{n}(x)\to\varphi(x)$, for each $x$, as $n\to\infty$, where $\varphi_{n}$ is the characteristic function corresponding to the distribution function $F_{n}$, and if $\varphi$ is continuous at the origin, then $F_{n}$ converges to a distribution function $F$, at the points of continuity of $F$, and $\varphi$ is the characteristic function of $F$.
\end{theorem}

\begin{proof}
Since $F_{n}(x)$ is a non-decreasing function of $x$ for $n=1,2,\dots$, and since $0\le F_{n}(x)\le1$, for all $x\in(-\infty,\infty)$, there exists (by Helly's theorem) a subsequence $F_{n_{k}}$ of $F_{n}$ of non-decreasing functions (of $x$) which converges, as $n_{k}\to\infty$, to a non-decreasing function $F(x)$. Clearly $0\le F(x)\le1$. We shall see that $F$ is a distribution function, and that $F_{n}\to F$, at every point of continuity of $F$.

By definition, we have
\[
\varphi_{n}(x)=\int_{-\infty}^{\infty}e^{ixy}\,dF_{n}(y).
\]
If we use formula \eqref{eq:III-1.8}, we get, for $R>0$,
\begin{equation}\tag{2.4}\label{eq:III-2.4}
\int_{0}^{R}\bigl\{F_{n_{k}}(y)-F_{n_{k}}(-y)\bigr\}\,dy
=\frac{1}{\pi}\int_{-\infty}^{\infty}\varphi_{n_{k}}(y)\,\frac{1-\cos Ry}{y^{2}}\,dy.
\end{equation}
Since $|\varphi_{n}(x)|\le1$, for all $x\in(-\infty,\infty)$, and $\varphi_{n}(x)\to\varphi(x)$, for each $x$, as $n\to\infty$, we get on letting $n_{k}\to\infty$ in \eqref{eq:III-2.4},
\[
\int_{0}^{R}\bigl[F(y)-F(-y)\bigr]\,dy=\frac{1}{\pi}\int_{-\infty}^{\infty}\varphi(y)\,\frac{1-\cos Ry}{y^{2}}\,dy,
\]
or
\begin{equation}\tag{2.5}\label{eq:III-2.5}
\begin{aligned}
&\frac{1}{R}  \int_{0}^{R}  \bigl[F(y)-F(-y)\bigr]\,dy\\
&= \frac{1}{2\pi}  \int_{-\infty}^{\infty}\varphi(y)\,H_{R}(y)\,dy,  \qquad H_{R}(y) = \frac{\sin^{2}(Ry/2)}{R(y/2)^{2}}.
\end{aligned}
\end{equation}
Since $\varphi_{n}(0)=1$ for $n=1,2,\dots$, and $\varphi(x)=\lim_{n\to\infty}\varphi_{n}(x)$, we have $\varphi(0)=1$, and since $\int_{-\infty}^{\infty}\bigl(\frac{\sin\alpha}{\alpha}\bigr)^{2}d\alpha=\pi$, we have $\frac{1}{2\pi}\int_{-\infty}^{\infty}H_{R}(y)\,dy=1$. (cf.\ Ex.1, Th.13, Ch.~I). Hence, for a suitably chosen $\delta>0$, we have
\begin{align*}
&\frac{1}{2\pi}\int_{-\infty}^{\infty}\varphi(y)\,H_{R}(y)\,dy-1
=\frac{1}{2\pi}\int_{-\infty}^{\infty}\bigl\{\varphi(y)-\varphi(0)\bigr\}\,H_{R}(y)\,dy\\
&=\frac{1}{2\pi}\Bigl(\int_{|y|<\delta}+\int_{|y|\ge\delta}\Bigr)\bigl\{\varphi(y)-\varphi(0)\bigr\}\,H_{R}(y)\,dy\\
&=I_{1}+I_{2},\quad\text{say.}
\end{align*}
Because of the continuity of $\varphi$ at the origin, given $\varepsilon>0$, there exists a $\delta>0$, such that
\[
|I_{1}|\le\sup_{|y|<\delta}\bigl|\varphi(y)-\varphi(0)\bigr|\cdot1<\varepsilon,
\]
while
\[
|I_{2}|\le\frac{1}{\pi}\int_{|t|\ge R\delta}\Bigl(\frac{\sin t/2}{t/2}\Bigr)^{2}dt\to0,
\]
as $R\to\infty$ for a suitably chosen, but fixed, $\delta>0$. Hence
\[
\lim_{R\to\infty}\frac{1}{2\pi}\int_{-\infty}^{\infty}\varphi(y)\,H_{R}(y)\,dy=1,
\]
which implies, because of \eqref{eq:III-2.5}, that
\begin{equation}\tag{2.6}\label{eq:III-2.6}
\lim_{R\to\infty}\frac{1}{R}\int_{0}^{R}\bigl\{F(y)-F(-y)\bigr\}\,dy=1.
\end{equation}
The left-hand side equals $F(+\infty)-F(-\infty)$. Since $0\le F(x)\le1$, and $F$ is non-decreasing, we deduce that $F(+\infty)=1$, $F(-\infty)=0$. Hence $F$ is a distribution function. By Theorem~\ref{thm:III-3}, we have
\[
\lim_{n_{k}\to\infty}\varphi_{n_{k}}(x)=\varphi(x)=\int_{-\infty}^{\infty}e^{ixy}\,dF(y).
\]
If there exists a point of continuity $x_{0}$ of $F$, such that $F_{n}(x_{0})\not\to F(x_{0})$, then we can find a subsequence $F_{n_{k}'}$, which converges everywhere to a distribution function $F^{*}$, and such that $F^{*}(x_{0})\neq F(x_{0})$. If, for example, $F^{*}(x_{0})>F(x_{0})$, then since $F$ is continuous at $x_{0}$, and $F^{*}$ is non-decreasing, we have $F^{*}(x)>F(x)$, for $x\in(x_{0},x_{0}+\eta)$, for some $\eta>0$, sufficiently small. (Similarly if $F^{*}(x_{0})<F(x_{0})$, then $F^{*}(x)<F(x)$ for $x\in(x_{0}-\eta',x_{0})$ for some $\eta'>0$). This is impossible since, by Theorem~\ref{thm:III-3}, $\varphi$ is also the characteristic function of $F^{*}$, and the characteristic function determines the distribution function up to an additive constant (Theorem~\ref{thm:III-1}) and $F^{*}(-\infty)=F(-\infty)=0$.
\end{proof}

\section{Positive definite functions: the theorems of Bochner and of F. Riesz}
\label{sec:III-3}

We shall consider classes of functions $f$ which can be represented as Fourier--Stieltjes transforms.

\begin{lemma}\label{lem:III-1}
If $f(t)$ is a complex-valued function which is measurable and finite for $-\infty<t<\infty$, and satisfies the condition
\begin{equation}\tag{3.1}\label{eq:III-3.1}
\sum_{\mu=1}^{m}\sum_{\nu=1}^{m}f(t_{\mu}-t_{\nu})\,\rho_{\mu}\,\overline{\rho}_{\nu}\ge0
\end{equation}
for any integer $m\ge1$, and arbitrary real numbers $t_{1},t_{2},\dots,t_{m}$ and arbitrary complex numbers $\rho_{1},\rho_{2},\dots,\rho_{m}$, then we have
\begin{equation}\tag{3.2}\label{eq:III-3.2}
f(0)\ge0;\qquad f(-t)=\overline{f(t)};\qquad |f(t)|\le f(0).
\end{equation}
\end{lemma}

\begin{proof}
On taking $m=1$ in \eqref{eq:III-3.1}, we obtain
\[
|\rho_{1}|^{2}f(0)\ge0,
\]
so that $f(0)\ge0$.

On taking $m=2$, $t_{1}=t$, $t_{2}=0$, we obtain
\begin{equation}\tag{3.3}\label{eq:III-3.3}
f(0)\bigl[|\rho_{1}|^{2}+|\rho_{2}|^{2}\bigr]+f(t)\rho_{1}\overline{\rho}_{2}+f(-t)\overline{\rho}_{1}\rho_{2}\ge0.
\end{equation}
If we choose $\rho_{2}=1$ in \eqref{eq:III-3.3}, we get
\[
f(0)\bigl[|\rho_{1}|^{2}+1\bigr]+f(t)\rho_{1}+f(-t)\overline{\rho}_{1}\ge0,
\]
and if in this we set $\rho_{1}=1$, $i$ respectively, we see that $f(t)+f(-t)$ is real (since $f(0)$ is real), and that $f(t)-f(-t)$ is purely imaginary, hence $f(-t)=\overline{f(t)}$.

If $f(0)=0$, $\rho_{1}=1$, $\rho_{2}=-f(t)$, then \eqref{eq:III-3.3} gives $-2|f(t)|^{2}\ge0$, since $f(-t)=\overline{f(t)}$, hence $f(t)=0$, for all $t$.

If $f(0)>0$, $\rho_{1}=f(0)$, $\rho_{2}=-f(t)$, then \eqref{eq:III-3.3} gives
\[
f(0)\bigl[|f(0)|^{2}+|f(t)|^{2}\bigr]>f(0)\bigl\{2|f(t)|\bigr\}^{2},
\]
which implies that $|f(t)|\le f(0)$, and \eqref{eq:III-3.2} is proved.
\end{proof}

\begin{lemma}[F. Riesz]\label{lem:III-2}
For any complex-valued measurable function $f$, condition \eqref{eq:III-3.1} implies that
\begin{equation}\tag{3.4}\label{eq:III-3.4}
\int_{-\infty}^{\infty}\int_{-\infty}^{\infty}f(t-s)\,\rho(t)\,\overline{\rho}(s)\,ds\,dt\ge0,
\end{equation}
for any $\rho\in L_{1}(-\infty,\infty)$, provided that $f(0)$ is finite.
\end{lemma}

\begin{proof}
If $A>0$, and $\rho\in L_{2}(-A,A)$, we take $\rho_{\mu}=\rho(t_{\mu})$ in \eqref{eq:III-3.1} and integrate with respect to $t_{\mu}$, $\mu=1,2,\dots,m$, where $m>1$. We get from each of the diagonal terms (in which $\mu=\nu$)
\[
(2A)^{m-1}f(0)\int_{-A}^{A}|\rho(t)|^{2}\,dt,
\]
while each of the remaining terms gives
\[
(2A)^{m-2}\int_{-A}^{A}\int_{-A}^{A}f(t-s)\,\rho(t)\,\overline{\rho}(s)\,ds\,dt.
\]
Hence we have
\[
\begin{aligned}
&m(2A)^{m-1}f(0)\int_{-A}^{A}|\rho(t)|^{2}\,dt\\
&\qquad+m(m-1)(2A)^{m-2}\int_{-A}^{A}\int_{-A}^{A}f(t-s)\,\rho(t)\,\overline{\rho}(s)\,ds\,dt\ge0.
\end{aligned}
\]
On dividing throughout by $m(m-1)(2A)^{m-2}$, and then letting $m\to\infty$, we get
\begin{equation}\tag{3.5}\label{eq:III-3.5}
\int_{-A}^{A}\int_{-A}^{A}f(t-s)\,\rho(t)\,\overline{\rho}(s)\,ds\,dt\ge0,
\end{equation}
for any $\rho\in L_{2}(-A,A)$. If $\rho\in L_{1}(-A,A)$, \eqref{eq:III-3.5} still holds good, for if we define, for any integer $n\ge1$,
\[
\rho_{n}(t)=
\begin{cases}
\rho(t),&\text{if }|\rho(t)|\le n,\\
n,&\text{if }|\rho(t)|>n,
\end{cases}
\]
then we have $|\rho_{n}(t)|\le|\rho(t)|\in L_{1}(-A,A)$, and $|\rho_{n}(t)|\le n$, for all $t\in(-A,A)$. Hence $\rho_{n}\in L_{2}(-A,A)\cap L_{1}(-A,A)$, so that \eqref{eq:III-3.5} holds with $\rho_{n}(t)$ in place of $\rho(t)$. Since $\rho_{n}(t)\to\rho(t)$, as $n\to\infty$, and $|f(t)|\le f(0)$, and $f(0)$ is finite, we obtain \eqref{eq:III-3.5} for any $\rho\in L_{1}(-A,A)$. If we then let $A\to\infty$, we obtain \eqref{eq:III-3.4}.
\end{proof}

\begin{lemma}\label{lem:III-3}
If a measurable function $f$ (as in Lemma~\ref{lem:III-1}) satisfies \eqref{eq:III-3.1}, and $f(0)$ is finite, then for any $\varepsilon>0$, the function
\begin{equation}\tag{3.6}\label{eq:III-3.6}
f_{\varepsilon}(x)=e^{-\varepsilon x^{2}}f(x),\qquad-\infty<x<\infty,\ \varepsilon>0,
\end{equation}
also satisfies \eqref{eq:III-3.1}; and $f_{\varepsilon}(x)\in L_{1}(-\infty<x<\infty)$.
\end{lemma}

\begin{proof}
We have
\begin{align*}
&\sum_{\mu=1}^{m}\sum_{\nu=1}^{m}f_{\varepsilon}(t_{\mu}-t_{\nu})\,\rho_{\mu}\,\overline{\rho}_{\nu}
=\sum_{\mu=1}^{m}\sum_{\nu=1}^{m}\rho_{\mu}\,\overline{\rho}_{\nu}\,f(t_{\mu}-t_{\nu})\,e^{-\varepsilon(t_{\mu}-t_{\nu})^{2}}\\
&=\sum_{\mu=1}^{m}\sum_{\nu=1}^{m}\rho_{\mu}\,\overline{\rho}_{\nu}\,f(t_{\mu}-t_{\nu})\,\\
&\qquad\frac{1}{2(\pi\varepsilon)^{1/2}}\int_{-\infty}^{\infty}e^{-t^{2}/4\varepsilon+i(t_{\mu}-t_{\nu})t}\,dt,
\end{align*}
if we note that
\begin{equation}\tag{3.7}\label{eq:III-3.7}
\int_{-\infty}^{\infty}e^{-x^{2}+i\alpha x}\,dx=\sqrt{\pi}\,e^{-\alpha^{2}/4},\qquad(\text{see }\S\ref{sec:I-1},\text{ Ch.~I})
\end{equation}
and replace $x$ by $x/(2\sqrt{\varepsilon})$, and $\alpha$ by $2\alpha\sqrt{\varepsilon}$.

Now the expression in \eqref{eq:III-3.7} equals the integral
\[
\int_{-\infty}^{\infty}e^{-t^{2}/4\varepsilon}
\Biggl(\sum_{\mu=1}^{m}\sum_{\nu=1}^{m}
\bigl(\rho_{\mu}e^{it_{\mu}t}\bigr)\overline{\bigl(\rho_{\nu}e^{it_{\nu}t}\bigr)}\,f(t_{\mu}-t_{\nu})\Biggr)\,dt\ge0,
\]
since $f$ satisfies \eqref{eq:III-3.1} by hypothesis. It follows from \eqref{eq:III-3.7} that $f_{\varepsilon}$ satisfies \eqref{eq:III-3.1}. Since $f(0)$ is finite, and (by Lemma~\ref{lem:III-1}) $|f(t)|\le f(0)$, we note that $f$ is bounded, hence $f_{\varepsilon}\in L_{1}(-\infty,\infty)$ for every $\varepsilon>0$.
\end{proof}

\begin{lemma}\label{lem:III-4}
Suppose that $f$ is a complex-valued measurable function on $(-\infty,\infty)$ which satisfies \eqref{eq:III-3.1}, with $f(0)$ finite. Let $f_{\varepsilon}(x)=e^{-\varepsilon x^{2}}f(x)$, $\varepsilon>0$, so that (by Lemma~\ref{lem:III-3}) $f_{\varepsilon}\in L_{1}(-\infty,\infty)$. Then we have
\[
\hat f_{\varepsilon}(\alpha)\ge0,\qquad \hat f_{\varepsilon}\in L_{1}(-\infty,\infty),\qquad-\infty<\alpha<\infty,
\]
and
\begin{equation}\tag{3.8}\label{eq:III-3.8}
f_{\varepsilon}(x)=\frac{1}{2\pi}\int_{-\infty}^{\infty}\hat f_{\varepsilon}(\alpha)\,e^{-i\alpha x}\,d\alpha,
\end{equation}
for almost all $x\in(-\infty,\infty)$. In particular, \eqref{eq:III-3.8} holds at every point of continuity $x$ of $f_{\varepsilon}$, and therefore of $f$.
\end{lemma}

\begin{proof}
By Lemma~\ref{lem:III-2}, $f$ satisfies condition \eqref{eq:III-3.4}. In it we choose $\rho$, such that
\[
\rho(t)=e^{-2\varepsilon t^{2}}e^{i\alpha t},\qquad-\infty<t<\infty,
\]
where $\alpha$ is real and fixed. Then \eqref{eq:III-3.4} becomes
\[
\int_{-\infty}^{\infty}\int_{-\infty}^{\infty}f(x-y)\,e^{-2\varepsilon(x^{2}+y^{2})}\,e^{i\alpha(x-y)}\,dx\,dy\ge0.
\]
If we make the substitutions
\[
x-y=u,\qquad x+y=v,
\]
then we have
\begin{align*}
&\frac12\int_{-\infty}^{\infty}\int_{-\infty}^{\infty}f(u)\,e^{-\varepsilon(u^{2}+v^{2})}\,e^{iu\alpha}\,du\,dv\\
&=\frac12\int_{-\infty}^{\infty}\Bigl(\int_{-\infty}^{\infty}e^{-\varepsilon v^{2}}\,dv\Bigr)e^{-\varepsilon u^{2}}f(u)\,e^{iu\alpha}\,du\\
&=\frac12\int_{-\infty}^{\infty}\Bigl(\int_{-\infty}^{\infty}e^{-\varepsilon v^{2}}\,dv\Bigr)f_{\varepsilon}(u)\,e^{iu\alpha}\,du\ge0.
\end{align*}
Now $f_{\varepsilon}\in L_{1}(-\infty,\infty)$, and $f_{\varepsilon}$ is bounded on $(-\infty,\infty)$, since $f(0)$ is finite, and the last inequality shows that $f_{\varepsilon}$ has a positive Fourier transform. By Theorem~12 of Ch.~I, $\hat f_{\varepsilon}\in L_{1}(-\infty,\infty)$ and \eqref{eq:III-3.8} follows.
\end{proof}

\begin{lemma}\label{lem:III-5}
Let
\[
K(\alpha)=
\begin{cases}
1-|\alpha|,&\text{for }|\alpha|\le1,\\
0,&\text{for }|\alpha|>1,
\end{cases}
\]
so that $K\in L_{1}(-\infty,\infty)$, and let $\hat K=H$. If, for $R>0$, $K_{R}(\alpha)\equiv K(\alpha/R)$, and $H_{R}(\alpha)=RH(R\alpha)$, then
\[
\hat K_{R}=H_{R}\qquad\text{and}\qquad\hat H_{R}=2\pi\,K_{R}.
\]
(If $f\in L_{1}(-\infty,\infty)$, $\hat f(\alpha)=\int_{-\infty}^{\infty}f(x)e^{i\alpha x}\,dx$, as in Ch.~I, \eqref{eq:I-8.1}). If $f_{\varepsilon}$ is defined for every $\varepsilon>0$ as in Lemma~\ref{lem:III-3}, then we have
\begin{equation}\tag{3.9}\label{eq:III-3.9}
\int_{-\infty}^{\infty}H_{R}(x-y)\,f_{\varepsilon}(y)\,dy
=\int_{-\infty}^{\infty}K_{R}(\alpha)\,e^{i\alpha x}\,\hat f_{\varepsilon}(-\alpha)\,d\alpha.
\end{equation}
\end{lemma}

\begin{proof}
Since $K$ is even, $H$ is even, so that the left-hand side of \eqref{eq:III-3.9} equals
\[
\int_{-\infty}^{\infty}H_{R}(y-x)\,f_{\varepsilon}(y)\,dy
=\int_{-\infty}^{\infty}f_{\varepsilon}(x+t)\,H_{R}(t)\,dt.
\]
By the composition rule (Ch.~I, \eqref{eq:I-1.13}), this equals
\[
\int_{-\infty}^{\infty}K_{R}(\alpha)\,e^{-i\alpha x}\,\hat f_{\varepsilon}(\alpha)\,d\alpha
=\int_{-\infty}^{\infty}K_{R}(\alpha)\,e^{+i\alpha x}\,\hat f_{\varepsilon}(-\alpha)\,d\alpha,
\]
since $f_{\varepsilon}\in L_{1}(-\infty,\infty)$ by Lemma~\ref{lem:III-4}, $K_{R}\in L_{1}(-\infty,\infty)$, $H_{R}\in L_{1}(-\infty,\infty)$, $K_{R}$ is even, and $\hat K_{R}=H_{R}$.
\end{proof}

\begin{lemma}\label{lem:III-6}
Let $f$ be a complex-valued measurable function defined on $(-\infty,\infty)$, which satisfies condition \eqref{eq:III-3.1}, and is continuous for $x=0$. Let
\begin{equation}\tag{3.10}\label{eq:III-3.10}
f_{n}(x)=e^{-x^{2}/n}f(x),\qquad n\ge1,\ n\text{ integral},\ -\infty<x<\infty.
\end{equation}
Then there exists a non-decreasing, bounded function $V_{n}(t)$, $-\infty<t<\infty$, such that
\begin{equation}\tag{3.11}\label{eq:III-3.11}
\begin{aligned}
&\frac{1}{2\pi}  \int_{-\infty}^{\infty}H_{R}(x-y)\,f_{n}(y)\,dy\\
&= \int_{-\infty}^{\infty}K_{R}(\alpha)\,e^{i\alpha x}\,dV_{n}(\alpha),  \qquad - \infty < x < \infty.
\end{aligned}
\end{equation}
\end{lemma}

\begin{proof}
If $f$ is continuous at $x=0$, then $f(0)$ is finite, and by Lemma~\ref{lem:III-4}, $\hat f_{n}(-\alpha)\ge0$, for $-\infty<\alpha<\infty$, $\hat f_{n}\in L_{1}(-\infty,\infty)$. If we define
\begin{equation}\tag{3.12}\label{eq:III-3.12}
V_{n}(t)=\frac{1}{2\pi}\int_{-\infty}^{t}\hat f_{n}(-\alpha)\,d\alpha,\qquad-\infty<t<\infty,
\end{equation}
then $V_{n}(t)\ge0$, and $V_{n}(t)$ is a non-decreasing function of $t$ for $n=1,2,\dots$. If $f$ is continuous at the origin, then, by Lemma~\ref{lem:III-4}, we have
\[
f_{n}(0)=\frac{1}{2\pi}\int_{-\infty}^{\infty}\hat f_{n}(-\alpha)\,d\alpha,
\]
so that
\begin{equation}\tag{3.13}\label{eq:III-3.13}
0\le V_{n}(t)\le f_{n}(0)=f(0),\qquad\text{for }-\infty<t<\infty,\ n=1,2,\dots.
\end{equation}
Now \eqref{eq:III-3.9} gives \eqref{eq:III-3.11}.
\end{proof}

\begin{theorem}\label{thm:III-5}
Let $f$ be a complex-valued measurable function defined on $(-\infty,\infty)$, which is continuous at the origin, and which satisfies condition \eqref{eq:III-3.1}. Then there exists a non-decreasing, bounded function $V(t)$, such that
\begin{equation}\tag{3.14}\label{eq:III-3.14}
f(x)=\int_{-\infty}^{\infty}e^{ixt}\,dV(t),\qquad-\infty<x<\infty,
\end{equation}
for almost all $x\in(-\infty,\infty)$. If $f$ is continuous everywhere, then \eqref{eq:III-3.14} holds for every $x\in(-\infty,\infty)$.
\end{theorem}

\begin{proof}
We have to let $n\to\infty$ in relation \eqref{eq:III-3.11} of Lemma~\ref{lem:III-6} just proved, and then let $R\to\infty$. Since $\{V_{n}(t)\}$ is uniformly bounded, by \eqref{eq:III-3.13}, by a theorem of Helly there exists a subsequence $(n_{k})$, and a non-decreasing function $V$, such that $|V(t)|\le f(0)$, and $V_{n_{k}}\to V$ pointwise on $(-\infty,\infty)$.

Since $K_{R}(\alpha)$ in Lemma~\ref{lem:III-5} vanishes for $|\alpha|\ge R$, and is continuous, we have, by a theorem of Helly--Bray (see the Notes),
\begin{equation}\tag{3.15}\label{eq:III-3.15}
\lim_{k\to\infty}\int_{-\infty}^{\infty}K_{R}(\alpha)\,e^{i\alpha x}\,dV_{n_{k}}(\alpha)
=\int_{-\infty}^{\infty}K_{R}(\alpha)\,e^{i\alpha x}\,dV(\alpha),
\end{equation}
for each $R>0$, on the right-hand side of \eqref{eq:III-3.11}. If we take the left-hand side of \eqref{eq:III-3.11}, and let $n\to\infty$, through the same subsequence $(n_{k})$ as in \eqref{eq:III-3.15}, we obtain, since $|f_{n}(y)|\le f_{n}(0)=f(0)$, and $f(0)$ is finite, and $H_{R}\in L_{1}(-\infty,\infty)$,
\begin{equation}\tag{3.16}\label{eq:III-3.16}
\lim_{k\to\infty}\frac{1}{2\pi}\int_{-\infty}^{\infty}H_{R}(x-y)\,f_{n_{k}}(y)\,dy
=\frac{1}{2\pi}\int_{-\infty}^{\infty}H_{R}(x-y)\,f(y)\,dy,
\end{equation}
by Lebesgue's theorem on dominated convergence. From \eqref{eq:III-3.16} and \eqref{eq:III-3.15}, and \eqref{eq:III-3.11}, we thus obtain
\begin{equation}\tag{3.17}\label{eq:III-3.17}
\frac{1}{2\pi}\int_{-\infty}^{\infty}H_{R}(x-y)\,f(y)\,dy
=\int_{-\infty}^{\infty}K_{R}(\alpha)\,e^{i\alpha x}\,dV(\alpha),\qquad(R>0).
\end{equation}
We now let $R\to\infty$. The function $K_{R}(\alpha)e^{i\alpha x}$ is continuous in $\alpha$, and vanishes outside $(-R,R)$, and $K_{R}(\alpha)\to1$, as $R\to\infty$, for each $\alpha$, while
\[
\int_{-\infty}^{\infty}\bigl|dV(\alpha)\bigr|\le V(+\infty)-V(-\infty)\le f(0)<\infty.
\]
Hence
\begin{equation}\tag{3.18}\label{eq:III-3.18}
\lim_{R\to\infty}\int_{-\infty}^{\infty}K_{R}(\alpha)\,e^{i\alpha x}\,dV(\alpha)
=\int_{-\infty}^{\infty}e^{i\alpha x}\,dV(\alpha),
\end{equation}
on the right-hand side of \eqref{eq:III-3.17}. We shall show that the left-hand side of \eqref{eq:III-3.17} gives
\begin{equation}\tag{3.19}\label{eq:III-3.19}
\lim_{R\to\infty}\frac{1}{2\pi}\int_{-\infty}^{\infty}f(y)\,H_{R}(x-y)\,dy
=\lim_{R\to\infty}\frac{1}{2\pi}(f*H_{R})(x)=f(x),
\end{equation}
for almost all $x\in(-\infty,\infty)$, so that \eqref{eq:III-3.14} follows, and the theorem is proved.

Let $\omega>0$, arbitrarily chosen and kept fixed, and let $f_{\omega}(y)=f(y)$, for $|y|<\omega$; while $f_{\omega}(y)=0$, for $|y|\ge\omega$. Since $f$ is bounded and measurable in $(-\omega,\omega)$, $f\in L_{1}(-\omega,\omega)$; hence $f_{\omega}\in L_{1}(-\infty,\infty)$. By Theorem~10 of Ch.~I, we have
\[
\lim_{R\to\infty}\frac{1}{2\pi}(f_{\omega}*H_{R})(x)=f_{\omega}(x),\qquad\text{for almost all }x\in(-\infty,\infty).
\]
That is to say that
\begin{equation}\tag{3.20}\label{eq:III-3.20}
\lim_{R\to\infty}\frac{1}{2\pi}\int_{-\omega}^{\omega}f(y)\,H_{R}(x-y)\,dy=f(x),
\end{equation}
for almost all $x\in[-\omega,\omega]$. By a further conclusion in Theorem~10 of Ch.~I, this relation holds for every $x\in(-\omega,\omega)$, if $f$ is everywhere continuous.

On the other hand, if $|x|<\omega$, we have
\begin{equation}\tag{3.21}\label{eq:III-3.21}
\begin{aligned}
&\Bigl|\frac{1}{2\pi}\int_{|y|\ge\omega}f(y)\,H_{R}(x-y)\,dy\Bigr|
\le\frac{f(0)}{2\pi\,R}\int_{|y|\ge\omega}\frac{\sin^{2}[R(x-y)/2]}{[(x-y)/2]^{2}}\,dy\to0,\\
&\qquad\text{as }R\to\infty.
\end{aligned}
\end{equation}
Thus \eqref{eq:III-3.21} and \eqref{eq:III-3.20} lead to the relation
\begin{equation}\tag{3.19}\label{eq:III-3.19b}
\lim_{R\to\infty}\frac{1}{2\pi}(f*H_{R})(x)=f(x),
\end{equation}
for almost all $x\in[-\omega,\omega]$, where $\omega>0$ is arbitrary; and for every $x\in(-\omega,\omega)$ if $f$ is continuous everywhere. This combined with \eqref{eq:III-3.18} and \eqref{eq:III-3.17} proves the theorem.
\end{proof}

The next result is a kind of converse and easier to prove.

\begin{theorem}[Bochner]\label{thm:III-6}
Let $V(t)$ be a non-decreasing function of $t$, $-\infty<t<\infty$, which is bounded everywhere. Then the function
\begin{equation}\tag{3.22}\label{eq:III-3.22}
f(x)=\int_{-\infty}^{\infty}e^{ixt}\,dV(t),\qquad-\infty<x<\infty,
\end{equation}
is continuous, and bounded, and satisfies condition \eqref{eq:III-3.1}.
\end{theorem}

\begin{proof}
As already remarked in \S\ref{sec:III-1}, $f$ is bounded and continuous. To verify \eqref{eq:III-3.1} we have only to note that
\begin{align*}
&\sum_{\mu=1}^{m}\sum_{\nu=1}^{m}f(t_{\mu}-t_{\nu})\,\rho_{\mu}\,\overline{\rho}_{\nu}
=\sum_{\mu=1}^{m}\sum_{\nu=1}^{m}\rho_{\mu}\,\overline{\rho}_{\nu}
\int_{-\infty}^{\infty}e^{it(t_{\mu}-t_{\nu})}\,dV(t)\\
&=\int_{-\infty}^{\infty}\Bigl(\sum_{\mu=1}^{m}\rho_{\mu}e^{it_{\mu}t}\Bigr)
\Bigl(\sum_{\nu=1}^{m}\overline{\rho}_{\nu}e^{-it_{\nu}t}\Bigr)\,dV(t)\\
&=\int_{-\infty}^{\infty}\Bigl|\sum_{\mu=1}^{m}\rho_{\mu}e^{it_{\mu}t}\Bigr|^{2}\,dV(t)\ge0,
\end{align*}
since $V$ is non-decreasing, and the integrand is positive.
\end{proof}

\begin{nstmt}{3.23}\label{eq:III-3.23}
Definition. A complex-valued function $f(x)$ defined for $-\infty<x<\infty$ is said to be \emph{positive definite}, if it is continuous, and bounded, and satisfies condition \eqref{eq:III-3.1}.
\end{nstmt}

Theorems \ref{thm:III-5} and \ref{thm:III-6} yield the following

\begin{theorem}[Bochner]\label{thm:III-7}
In order that a function $f(x)$ defined on $(-\infty,\infty)$ may be written as
\[
f(x)=\int_{-\infty}^{\infty}e^{ixt}\,dV(t),
\]
where $V$ is a non-decreasing, bounded function on $(-\infty,\infty)$, it is necessary and sufficient that $f$ be positive definite.
\end{theorem}

Since a positive definite function is, by definition, continuous in $(-\infty,\infty)$, the proof of Theorem~\ref{thm:III-7} does not require the full force of Theorem~\ref{thm:III-5}. The following theorem is sufficient.

\begin{theorem}\label{thm:III-8}
If $f$ is a complex-valued function defined on $(-\infty,\infty)$, which is continuous, and satisfies condition \eqref{eq:III-3.1}, then
\begin{equation}\tag{3.24}\label{eq:III-3.24}
f(x)=\int_{-\infty}^{\infty}e^{ixt}\,dV(t),
\end{equation}
where $V$ is a non-decreasing, bounded function of $t\in(-\infty,\infty)$.
\end{theorem}

\begin{proof}
Define for each integer $n$, $n>1$,
\[
f_{n}(x)=e^{-x^{2}/n}f(x),\qquad-\infty<x<\infty.
\]
Then we have $f_{n}\in L_{1}(-\infty,\infty)$, for every $n$, since $f$ is bounded, and by Lemma~\ref{lem:III-4}, with $\varepsilon=1/n$, we have $\hat f_{n}\in L_{1}(-\infty,\infty)$, and
\begin{equation}\tag{3.25}\label{eq:III-3.25}
f_{n}(x)=\frac{1}{2\pi}\int_{-\infty}^{\infty}\hat f_{n}(\alpha)\,e^{-i\alpha x}\,d\alpha
=\frac{1}{2\pi}\int_{-\infty}^{\infty}\hat f_{n}(-\alpha)\,e^{i\alpha x}\,d\alpha,
\end{equation}
where $\hat f_{n}(-\alpha)\ge0$, for $-\infty<\alpha<\infty$, for every $x\in(-\infty,\infty)$, since $f$ is continuous (see Th.11$'$, Ch.~I, p.~45).

If we define, as in \eqref{eq:III-3.12},
\[
V_{n}(t)=\frac{1}{2\pi}\int_{-\infty}^{t}\hat f_{n}(-\alpha)\,d\alpha,\qquad-\infty<t<\infty,
\]
then \eqref{eq:III-3.25} becomes
\begin{equation}\tag{3.26}\label{eq:III-3.26}
f_{n}(x)=\int_{-\infty}^{\infty}e^{i\alpha x}\,dV_{n}(\alpha).
\end{equation}
Now $V_{n}(-\infty)=0$, while $V_{n}(+\infty)=f_{n}(0)=f(0)$.

We consider two cases.

\noindent\textbf{Case (i).} Let $f(0)=1$. Then $V_{n}$ is a distribution function, for each $n$, with $f_{n}$ as its characteristic function. Since $f_{n}(x)\to f(x)$, as $n\to\infty$, for each $x$, and $f$ is continuous (at the origin), Theorem~\ref{thm:III-4} implies that $V_{n}$ converges to a distribution function $V$ (at the points of continuity of $V$), whose characteristic function is $f$. That is to say, \eqref{eq:III-3.26} implies that
\[
f(x)=\int_{-\infty}^{\infty}e^{i\alpha x}\,dV(\alpha)
\]
as claimed.

\noindent\textbf{Case (ii).} Let $f(0)\neq1$. Condition \eqref{eq:III-3.1} implies that $f(0)\ge0$. If $f(0)=0$, then $f$ is identically zero, and $V$ is a constant. If $f(0)\neq0$, then $f(0)>0$, and the function $g(x)=\frac{f(x)}{f(0)}$ is continuous, and satisfies condition \eqref{eq:III-3.1}, with $g(0)=1$. By what has been proved in Case~(i), it follows that
\[
f(x)=\int_{-\infty}^{\infty}e^{i\alpha x}\,d\bigl(f(0)\cdot V(\alpha)\bigr)
=\int_{-\infty}^{\infty}e^{i\alpha x}\,dV_{1}(\alpha),
\]
where $V_{1}(\alpha)=f(0)\cdot V(\alpha)$, and $V_{1}$ is non-decreasing, and bounded.
\end{proof}

The above proof yields the following

\begin{nstmt}{3.27}\label{eq:III-3.27}
Corollary. Positive definite functions $f$, with the property $f(0)=1$, are characteristic functions.
\end{nstmt}

An easy extension of Theorem~\ref{thm:III-7} is the following

\begin{theorem}[Bochner]\label{thm:III-9}
Let $g\in L_{2}(-\infty,\infty)$, and let
\begin{equation}\tag{3.28}\label{eq:III-3.28}
f(x)=\int_{-\infty}^{\infty}g(x+t)\,\overline{g(t)}\,dt.
\end{equation}
Then we have
\[
f(x)=\int_{-\infty}^{\infty}e^{ix\alpha}\,dV(\alpha),
\]
where $V$ is non-decreasing, and bounded, in $(-\infty,\infty)$.
\end{theorem}

\begin{proof}
Since
\[
|f(x)|^{2}\le\int_{-\infty}^{\infty}|g(x+t)|^{2}\,dt\int_{-\infty}^{\infty}|g(t)|^{2}\,dt=(f(0))^{2},
\]
which is finite, since $\|g\|_{2}$ is finite, $f$ is bounded.

We also have, for any fixed $x$,
\[
|f(x)-f(x_{1})|^{2}\le\int_{-\infty}^{\infty}|g(x+t)-g(x_{1}+t)|^{2}\,dt\int_{-\infty}^{\infty}|g(t)|^{2}\,dt,
\]
and if we let $x_{1}\to x$, then the right-hand side tends to zero, so that $f$ is continuous.

Since
\[
f(x-y)=\int_{-\infty}^{\infty}g(x-y+t)\,\overline{g}(t)\,dt
=\int_{-\infty}^{\infty}g(x+t)\,\overline{g(y+t)}\,dt,
\]
we have
\[
\sum_{\mu=1}^{m}\sum_{\nu=1}^{m}f(x_{\mu}-x_{\nu})\,\sigma_{\mu}\,\overline{\sigma}_{\nu}
=\int_{-\infty}^{\infty}\Bigl|\sum_{\mu=1}^{m}g(x_{\mu}+t)\sigma_{\mu}\Bigr|^{2}\,dt\ge0.
\]
It follows that $f$ is positive definite, and hence, by Theorem~\ref{thm:III-7}, that
\[
f(x)=\int_{-\infty}^{\infty}e^{ix\alpha}\,dV(\alpha).
\]
\end{proof}

\section{A uniqueness theorem}
\label{sec:III-4}

For a special kind of Fourier transform, namely
\[
\int_{-\infty}^{\infty}e^{ixy}\varphi(y)\,dy,
\]
where $\varphi\in L_{1}(-\omega,\omega)$, for each finite $\omega>0$, we shall prove a theorem of uniqueness, which generalizes Theorem~7 of Ch.~I.

\begin{theorem}[Offord]\label{thm:III-10}
If $\varphi(u)\in L_{1}(-\omega<u<\omega)$, for every finite $\omega>0$, and if
\begin{equation}\tag{4.1}\label{eq:III-4.1}
\lim_{\omega\to\infty}\int_{-\omega}^{\omega}\varphi(u)\,e^{iux}\,du=0,
\end{equation}
for every real $x$, then $\varphi(u)=0$, for almost every $u\in(-\infty,\infty)$.
\end{theorem}

We first prove three preliminary lemmas.

\begin{lemma}[Schwarz]\label{lem:III-7}
If $f(x)$ is real-valued and continuous for $a\le x\le b$, $b>a$, and
\begin{equation}\tag{4.2}\label{eq:III-4.2}
\lim_{h\to0}\frac{f(x+h)+f(x-h)-2f(x)}{h^{2}}=0,
\end{equation}
for every $x\in[a,b]$, then $f$ is linear in $(a,b)$; that is to say, $f(x)=Ax+B$, where $A$ and $B$ are some constants.
\end{lemma}

\begin{proof}
Let
\[
F(x)=\theta\bigl[f(x)-f(a)-\tfrac{x-a}{b-a}\bigl\{f(b)-f(a)\bigr\}\bigr],
\]
where $\theta=\pm1$, and let
\[
G(x)=F(x)-\tfrac12\varepsilon(x-a)(b-x),
\]
where $\varepsilon>0$. Then $G$ is a continuous function of $x$ for $x\in[a,b]$, with $G(a)=G(b)=0$; so is $F$, with $F(a)=F(b)=0$.

\noindent\textbf{Case (i).} If $F(x)=0$, for every $x\in[a,b]$, the lemma is immediate since then we have
\[
f(x)=\frac{f(b)-f(a)}{b-a}\cdot x+\frac{bf(a)-af(b)}{b-a}.
\]

\noindent\textbf{Case (ii).} If there exists a point $c$ in $(a,b)$, such that $F(c)\neq0$, choose $\theta$ in such a way that $F(c)>0$, and choose $\varepsilon$ so small that $G(c)>0$. Since $G$ is continuous in $[a,b]$, it attains its maximum $M$, say, at a point $x_{0}\in(a,b)$, and $M>0$, since $G(c)>0$. We have, because of \eqref{eq:III-4.2},
\[
\lim_{h\to0}\frac{G(x_{0}+h)+G(x_{0}-h)-2G(x_{0})}{h^{2}}=\varepsilon>0.
\]
But $G(x_{0}+h)\le G(x_{0})$, and $G(x_{0}-h)\le G(x_{0})$, so that the above limit is negative or zero, which implies a contradiction. Hence $F(x)=0$ for all $x\in[a,b]$, and the lemma follows.
\end{proof}

\begin{lemma}\label{lem:III-8}
Let $\varphi(u)\in L_{1}(-\omega<u<\omega)$, for each finite $\omega>0$, and let
\begin{equation}\tag{4.3}\label{eq:III-4.3}
\lim_{\omega\to\infty}I(\omega)\equiv\lim_{\omega\to\infty}\int_{-\omega}^{\omega}\varphi(u)\,du=0.
\end{equation}
If $\varphi(u)\to0$, as $|u|\to\infty$, then
\begin{equation}\tag{4.4}\label{eq:III-4.4}
\int_{-\infty}^{\infty}\varphi(u)\Bigl(\frac{\sin uh}{uh}\Bigr)^{2}du
\end{equation}
exists for every $h>0$, and tends to zero as $h\downarrow0$.
\end{lemma}

\begin{proof}
Since $\int_{-\omega}^{\omega}\varphi(u)\,du=\int_{0}^{\omega}\Phi(u)\,du$, where $\Phi(u)=\varphi(u)+\varphi(-u)$, assumption \eqref{eq:III-4.3} implies that $I(\omega)=\int_{0}^{\omega}\Phi(u)\,du\to0$, as $\omega\to\infty$. The integral in \eqref{eq:III-4.4} equals
\[
\int_{0}^{\infty}\Bigl(\frac{\sin uh}{uh}\Bigr)^{2}\,dI(u),
\]
and this, in turn, equals, by partial integration,
\[
\begin{aligned}
&\Bigl[I(u)\Bigl(\frac{\sin uh}{uh}\Bigr)^{2}\Bigr]_{u=0}^{\infty} - \int_{0}^{\infty}I(u)\,\frac{d}{du}  \Bigl[\Bigl(\frac{\sin uh}{uh}\Bigr)^{2}\Bigr]\,du\\
&= -  \int_{0}^{\infty}I  \Bigl(\frac{u}{h}\Bigr)\,H'(u)\,du,
\end{aligned}
\]
where $H'(u)\in L_{1}(0,\infty)$, $H(u)=\bigl(\frac{\sin u}{u}\bigr)^{2}$.

Since $I(u/h)$ is bounded, we have by Lebesgue's theorem on dominated convergence,
\[
\lim_{h\downarrow0}\int_{0}^{\infty}I\Bigl(\frac{u}{h}\Bigr)\,H'(u)\,du
=\int_{0}^{\infty}\lim_{h\downarrow0}I\Bigl(\frac{u}{h}\Bigr)\,H'(u)\,du=0.
\]
\end{proof}

\begin{lemma}\label{lem:III-9}
Let $\varphi(u)\in L_{1}(-\omega<u<\omega)$, for each finite $\omega>0$; let
\[
\lim_{\omega\to\infty}\int_{-\omega}^{\omega}\varphi(u)\,e^{iux}\,du=0,
\]
for every real $x$, and $\varphi(u)\to0$, as $|u|\to\infty$. Then $\varphi(u)=0$ for all $u\in(-\infty,\infty)$.
\end{lemma}

\begin{proof}
We take, as we may, $\varphi(u)e^{iux}$ instead of $\varphi(x)$ in Lemma~\ref{lem:III-8}. Then we have
\begin{equation}\tag{4.5}\label{eq:III-4.5}
\int_{-\infty}^{\infty}\varphi(u)\,e^{iux}\Bigl(\frac{\sin uh}{uh}\Bigr)^{2}du\to0,\qquad\text{as }h\downarrow0.
\end{equation}
Define
\[
F(x)=\int_{-\infty}^{\infty}\varphi(u)\,\frac{e^{iux}-L(x,u)}{-u^{2}}\,du,
\]
where
\[
L(x,u)=
\begin{cases}
1+iux,&\text{for }|u|<1,\\
0,&\text{for }|u|\ge1.
\end{cases}
\]
For a fixed $h\neq0$, consider the function
\begin{align*}
&\frac{F(x+h)+F(x-h)-2F(x)}{h^{2}}\\
&=\int_{-\infty}^{\infty}\frac{\varphi(u)}{(iuh)^{2}}\bigl\{e^{i(x+h)u}+e^{i(x-h)u}-2e^{iux}\bigr\}\,du\\
&=\int_{-\infty}^{\infty}\frac{\varphi(u)}{(iuh)^{2}}e^{iux}\Bigl\{e^{iuh/2}-e^{-iuh/2}\Bigr\}^{2}\,du\\
&=\int_{-\infty}^{\infty}\varphi(u)\,e^{iux}\Bigl(\frac{\sin(uh/2)}{uh/2}\Bigr)^{2}\,du.
\end{align*}
Because of \eqref{eq:III-4.5}, and Lemma~\ref{lem:III-7}, it follows that $F$ is linear, hence $F(x+h)+F(x-h)-2F(x)\equiv0$, which implies, in turn, that
\[
\int_{-\infty}^{\infty}\varphi(u)\,e^{iux}\Bigl(\frac{\sin\frac{uh}{2}}{uh/2}\Bigr)^{2}du\equiv0.
\]
But the integral on the left-hand side is the Fourier transform of a function belonging to $L_{1}(-\infty,\infty)$. Hence by Theorem~7 of Ch.~I, $\varphi(u)$ is zero for almost all $u\in(-\infty,\infty)$.
\end{proof}

We remark that this method of proof goes back to Riemann in his treatment of the uniqueness of trigonometric series.

\noindent\textbf{Proof of Theorem~\ref{thm:III-10}.}

Assumption \eqref{eq:III-4.1} holds with $y+x$, and $y-x$, in place of $x$ (where $y$ is real), so that it implies that
\begin{equation}\tag{4.7}\label{eq:III-4.7}
\int_{-\omega}^{\omega}\varphi(u)\,e^{iuy}\cos ux\,du\to0,\qquad\text{as }\omega\to\infty,
\end{equation}
for every real $y$ and $x$. If we set, for fixed $y$,
\[
\Psi(u)=\varphi(u)e^{iuy}+\varphi(-u)e^{-iuy},
\]
then \eqref{eq:III-4.7} becomes
\begin{equation}\tag{4.8}\label{eq:III-4.8}
\int_{0}^{\omega}\Psi(u)\cos ux\,du\to0,\qquad\text{as }\omega\to\infty,
\end{equation}
for every real $x$. Now let
\[
\Psi_{1}(v)=\int_{0}^{v}\Psi(u)\,du=\int_{-v}^{v}\varphi(u)\,e^{iuy}\,du,\qquad v>0.
\]
Then \eqref{eq:III-4.8} implies, in particular, that
\begin{equation}\tag{4.9}\label{eq:III-4.9}
\Psi_{1}(v)\to0,\qquad\text{as }v\to\infty.
\end{equation}
Since
\[
\int_{0}^{\omega}\Psi(u)\cos ux\,du
=\Psi_{1}(\omega)\cos\omega x-\Psi_{1}(0)+x\int_{0}^{\omega}\Psi_{1}(u)\sin ux\,du,
\]
where $\Psi_{1}(0)=0$, \eqref{eq:III-4.9} implies that
\[
x\int_{0}^{\omega}\Psi_{1}(u)\sin ux\,du\to0,\qquad\text{as }\omega\to\infty,
\]
which implies that for every $x\neq0$, we have
\[
\int_{0}^{\omega}\Psi_{1}(u)\sin ux\,du\to0,\qquad\text{as }\omega\to\infty,
\]
where $\Psi_{1}(u)\to0$, as $u\to\infty$, by \eqref{eq:III-4.9}. By Lemma~\ref{lem:III-9}, $\Psi_{1}(u)=0$, for almost all $u\in(-\infty,\infty)$. Since $\Psi_{1}$ is absolutely continuous, it follows that $\Psi(u)$ is zero for almost all $u$, hence $\varphi(u)e^{iuy}+\varphi(-u)e^{-iuy}=0$, for almost all $u$. On setting $y=0$, we have: $\varphi(u)+\varphi(-u)=0$, or $\varphi(u)=-\varphi(-u)$, for almost all $u$. Hence
\[
\Psi(u)=2i\,\varphi(u)\sin uy,
\]
and therefore $\varphi(u)=0$ for almost all $u\in(-\infty,\infty)$.

\begin{remark}\label{rem:III-1}
If $\varphi(u)=e^{iu^{2}}$, then $\varphi(u)\notin L_{1}(-\infty<u<\infty)$, although
\[
\lim_{\omega\to\infty}\int_{-\omega}^{\omega}\varphi(u)\,e^{iux}\,du
\]
is finite for all $x\in(-\infty,\infty)$. (See Ex.10, Step (iii), \S\ref{sec:I-1}, Ch.~I).

Another such example is provided by
\[
\varphi(u)=\exp(\alpha u+ie^{u}),\qquad0<\alpha<1.
\]
\end{remark}
